
\documentclass[12pt,onecolumn]{article}

\usepackage[utf8]{inputenc}
\usepackage[T1]{fontenc}
\usepackage{lmodern}
\usepackage[margin=1in]{geometry}
\usepackage{setspace}
\usepackage{microtype}
\usepackage{amsmath,amssymb,amsthm}
\usepackage{hyperref}
\usepackage{fancyhdr}
\usepackage{enumitem}
\usepackage{orcidlink}

\hypersetup{
    colorlinks=true,
    linkcolor=blue,
    citecolor=blue,
    urlcolor=blue
}

\setstretch{1.08}

\newtheorem{proposition}{Proposition}[section]
\newtheorem{corollary}[proposition]{Corollary}
\newtheorem{definition}[proposition]{Definition}
\newtheorem{remark}[proposition]{Remark}

\newcommand{\doi}[1]{\href{https://doi.org/#1}{#1}}

\pagestyle{fancy}
\fancyhf{}
\setlength{\headheight}{14pt}
\lhead{Does the Universe Have a Past?}
\rhead{John C. W. McKinley}
\cfoot{\thepage}

\title{
\textbf{Does the Universe Have a Past?}\\[0.5em]
\large Local Proper Time, Relativistic Registration,\\
and the Absence of a Physical Past
}

\author{
John C. W. McKinley\,
\orcidlink{0009-0005-7097-5035}
}

\date{October 2026}

\begin{document}

\maketitle

\begin{abstract}

The universe is ordinarily imagined as progressing along a timeline, with earlier physical configurations constituting its past, its current configuration constituting the present, and subsequent configurations constituting the future. This familiar picture encourages the assumption that the past continues to exist somewhere, as though earlier moments occupied locations along a physical timeline. The present paper challenges that assumption and proposes that the universe has no physically existing past.

The argument begins with an established result of relativity: proper time accumulates along individual timelike worldlines, not according to a universal clock. Two clocks following different spacetime paths between the same events can accumulate different durations. Special relativity also supplies no observer-independent simultaneity for spatially separated events, and general relativity supplies no universally preferred global present. The universe therefore possesses neither a universal accumulation of proper time nor a privileged sequence of instantaneous global configurations.

The physical significance of this result becomes clearer when applied locally. At every physical location, the registered configuration reflects the lawful history of physical interactions associated with that location. Time-bearing systems accumulate proper time along their respective worldlines, while the invariant limit \(c\) governs which causal information becomes available for local registration. Different locations consequently possess different causal relationships, temporal circumstances, and available information. There is no single universal present advancing through the cosmos. Under the proposed model, physical actuality consists of locally instantiated conditions and their lawful relationships.

Physics observes registered physical conditions, reconstructs earlier configurations from surviving evidence, and predicts possible subsequent configurations according to governing laws. Photographs, fossils, memories, and astronomical observations are present physical records, not the earlier configurations they describe. A mathematical timeline represents relationships among inferred and predicted configurations; it does not establish that those configurations exist as additional physical realities. The paper develops no-go arguments against a universal proper-time clock, a privileged global sequence of presents, and the inference that successful historical reconstruction establishes a physically existing past.

The result also clarifies the null proper-time foundation of the Timeless Light Model. There is no universal clock from which the photon needs an exemption. Emission and absorption possess registered spacetime relationships, but the photon relation accumulates no internal proper time. The apparent paradox arises from assigning a universal temporal progression to physical reality and then treating light as an exception. Removing that assumption eliminates the need for the exception.

The proposed architecture is a universe of locally registered physical conditions governed by timeless lawful structure. Earlier configurations are reconstructed, and subsequent configurations are predicted. Neither is an additional layer of present physical actuality.

The past is not somewhere else. It is not somewhere at all.

\end{abstract}

\noindent\textbf{Keywords:} proper time; simultaneity; local registration; physical history; historical reconstruction; spacetime; causal structure; ontology of time; null proper time; Timeless Light Model.

\section{Introduction: The Imaginary Timeline}

Human beings ordinarily imagine the universe as proceeding through time. Yesterday existed, today exists, and tomorrow will exist. In this familiar picture, the universe possesses a complete physical configuration at each moment, and those configurations are arranged along a timeline extending backward and forward from the observer's present position.

The picture is remarkably useful. A historian describes events in chronological order. An astronomer calculates the position of a planet at an earlier date. A physicist models a system through a sequence of states. A computer simulation represents successive configurations at successive values of a time parameter. These are successful applications of mathematical and historical description.

A difficulty arises when the representation is mistaken for the physical reality it represents.

Imagine a motion-picture film. Each frame records a complete visual configuration. The frames are arranged in a sequence, and a projector presents them one after another. The film remains available after projection. An earlier frame can be recovered, examined, and projected again.

The ordinary mental picture of time resembles this arrangement. The present is the frame currently being projected. The past consists of earlier frames. The future consists of later frames. The observer imagines moving along the film, or imagines the film moving past the observer.

But what establishes that the universe possesses such a film?

A physicist can calculate the Earth's position one hundred years ago. An astronomer can reconstruct conditions in the early observable universe. A historian can identify the circumstances surrounding an ancient event. None of these achievements establishes that the earlier physical configuration remains in existence at some location called the past.

The mathematical ability to identify an earlier configuration is not the physical ability to visit that configuration. The calculation describes a relationship among physical conditions. It does not establish that the universe has preserved the condition as an existing object.

Time-travel stories make the underlying assumption particularly visible. The traveler selects a year and journeys to it, as though the year were a destination with a continuing physical existence. The year 1850 is treated almost like a geographical location. The only missing element is the means of transportation.

Yet a date is not a physical location, and a coordinate assigned to an event is not evidence that the event remains physically available.

The present paper proposes removing that assumption from the interpretation of physics.

This is not an argument against historical science. The past is reconstructed with extraordinary success. Nor is it an argument against the mathematical use of timelines, spacetime coordinates, or differential equations containing time variables. Such descriptions remain indispensable.

The question is what physical existence those descriptions establish.

The argument begins from established relativistic physics rather than from a proposed metaphysical picture. Relativity establishes that proper time accumulates along timelike worldlines, that accumulated durations depend on the paths followed, and that no frame-independent universal simultaneity is supplied for spatially separated events. These results remove the physical basis for treating the universe as a collection of configurations advancing together under one universal clock.

The argument then turns to physical registration. Every observer receives information at a definite spacetime location. The information available there depends upon the causal circumstances of that location and upon the records preserved by physical systems. From these presently available conditions, physics reconstructs earlier events and predicts subsequent outcomes.

The central question follows naturally: what licenses the additional claim that the earlier configurations continue to exist?

The present paper develops a no-go result on that inference. It then advances a positive physical interpretation: the universe possesses locally actual physical conditions and the lawful structure governing them. Its past is not a physical region that continues to exist. Its future is not an already completed collection of configurations awaiting arrival.

The universe does not move along a stored timeline. The timeline is our description of lawful physical relationships.

\section{Proper Time Is Not Universal Time}

The most direct physical reason to reject a universal clock is that relativity does not assign one accumulated duration to every physical system.

In Newtonian mechanics, time is ordinarily represented by a parameter shared by the entire system. One writes the positions and velocities of objects as functions of a common variable \(t\). The resulting picture encourages the belief that the universe possesses one clock, even if individual observers employ different instruments to measure it.

Special relativity removes that foundation.

The physically accumulated duration of an ideal clock is its proper time. Proper time is measured along the clock's timelike worldline. It is not assigned independently of the clock's trajectory through spacetime.

For an ideal clock moving with speed \(v(t)\) relative to an inertial coordinate system in flat spacetime, the accumulated proper time is

\begin{equation}
\Delta\tau
=
\int_{t_1}^{t_2}
\sqrt{1-\frac{v(t)^2}{c^2}}\,dt.
\end{equation}

The coordinate time \(t\) provides a reference for describing the motion. The proper time \(\tau\) is the duration accumulated by the clock along its actual path.

This distinction is not a philosophical embellishment of relativity. It is part of the mathematical and experimentally supported structure of the theory.

Consider two ideal clocks initially together at the same spacetime event. Clock \(A\) remains at rest in an inertial frame. Clock \(B\) travels outward, reverses direction, and returns to meet clock \(A\). The clocks are compared when they reunite.

In the idealized case in which clock \(B\) travels at constant speed \(v\) during its outward and return segments, with a negligibly short turnaround, and the reunion occurs at inertial coordinate time \(2T\), their accumulated proper times are

\begin{equation}
\Delta\tau_A=2T
\end{equation}

and

\begin{equation}
\Delta\tau_B
=
2T\sqrt{1-\frac{v^2}{c^2}}.
\end{equation}

For nonzero \(v\),

\begin{equation}
\Delta\tau_A\neq\Delta\tau_B.
\end{equation}

The clocks depart from the same event and reunite at the same event. They nevertheless accumulate different amounts of proper time.

The significance of this result is easily underestimated because human beings are accustomed to interpreting clocks as imperfect instruments measuring a common underlying quantity. In this case, however, the difference does not arise from a mechanical defect or an inaccurate instrument. The clocks correctly accumulate different proper times because they follow different spacetime paths.

General relativity extends this description to curved spacetime. With metric signature \((-+++)\), the proper time accumulated along a timelike worldline \(\gamma\) is

\begin{equation}
\Delta\tau_\gamma
=
\frac{1}{c}
\int_\gamma
\sqrt{
-g_{\mu\nu}\,dx^\mu dx^\nu
}.
\end{equation}

The accumulated duration depends upon the metric and the worldline. Motion and gravitational circumstances both enter the calculation.

There is no requirement that two physically separated clocks accumulate identical durations. Their readings are governed by their own spacetime relationships.

This establishes the first result.

\begin{proposition}[No universal proper-time accumulation]

Relativistic physics does not supply a single accumulated proper time shared by all time-bearing physical systems.

\end{proposition}

\begin{proof}

Proper time is defined by integration along a timelike worldline. Different timelike worldlines connecting the same pair of events can possess different proper-time lengths, as demonstrated by the reunified-clock construction. Therefore accumulated proper time cannot be identified with one path-independent physical duration shared by every clock. The relativistic description requires individual worldline durations rather than one universally accumulated proper time.

\end{proof}

The proposition does not deny the usefulness of common coordinate time. A coordinate system can assign temporal labels to an extended spacetime region, and a particular cosmological model can employ a useful cosmic-time parameter. Such constructions organize the mathematical description. They do not make the proper times of all physical clocks identical.

The important physical distinction is between a coordinate convention and an accumulated physical duration.

Once that distinction is recognized, the universe no longer resembles a collection of objects all advancing through one shared internal clock.

Each time-bearing system accumulates its own proper time along its own worldline. The physical universe has no universal proper-time accumulation that every system must obey.

\section{The Positive No-Go on a Universal Clock}

The preceding result provides more than an argument from the absence of evidence. It identifies a positive physical fact that contradicts the familiar universal-clock picture.

If the universe possessed one physically accumulated time shared by all of its time-bearing constituents, clocks following different paths between the same two events would be required to accumulate that same universal duration. Relativity instead assigns them different proper times.

The physical structure is therefore not one clock accompanied by imperfect local measurements. It is a collection of time-bearing systems whose elapsed durations are determined by their individual spacetime trajectories.

This result changes the natural starting point for discussions of time.

The question should not be why an unusual physical entity fails to participate in universal temporal progression. No such universally accumulated temporal progression has been established.

The question should instead be which physical relations possess proper-time accumulation, which relations possess other forms of temporal ordering, and which relations possess no internal proper-time duration.

The familiar photon case will be considered later. Its significance becomes clearer after the universal-clock assumption has been removed.

For the present, the conclusion is straightforward.

\begin{corollary}[No universal clock is required]

The relativistic organization of physical events requires no single clock whose accumulated duration is shared by every physical system.

\end{corollary}

\begin{proof}

The governing relativistic relationships are expressed through spacetime geometry, causal structure, coordinate transformations, and proper times associated with particular worldlines. By the preceding proposition, those proper times need not agree. A universal proper-time clock therefore performs no necessary role in the description.

\end{proof}

This result concerns accumulated proper time. A second and complementary relativistic result concerns the meaning of a universal present.

Even if physical systems do not accumulate the same duration, one might still imagine that the entire universe occupies one objectively shared instant, with different clocks simply recording different elapsed times.

Special relativity removes that assumption as well.

\section{Relativity Does Not Supply a Universal Present}

A universal present would be one physical condition of simultaneity shared by the entire universe. It would identify which events at different locations belong to the same objectively privileged moment.

In special relativity, spatially separated simultaneity is not invariant under changes of inertial reference frame.

For two events separated by coordinate differences \(\Delta t\) and \(\Delta x\) in an inertial frame \(S\), the corresponding temporal separation in another inertial frame \(S'\), moving at speed \(v\) along the \(x\)-axis, is

\begin{equation}
\Delta t'
=
\gamma
\left(
\Delta t-\frac{v\Delta x}{c^2}
\right),
\end{equation}

where

\begin{equation}
\gamma
=
\frac{1}{\sqrt{1-v^2/c^2}}.
\end{equation}

Suppose the events are simultaneous in \(S\), so that \(\Delta t=0\). Their temporal separation in \(S'\) becomes

\begin{equation}
\Delta t'
=
-\gamma\frac{v\Delta x}{c^2}.
\end{equation}

For spatially separated events and nonzero relative velocity, this quantity is not zero.

Thus two spatially separated events assigned the same time coordinate in one inertial frame need not receive the same temporal assignment in another.

The theory supplies invariant spacetime intervals and causal relationships. It does not supply a single frame-independent assignment of distant simultaneity.

General relativity preserves the distinction. Spacetime geometry determines local causal relationships, light-cone structure, and proper-time accumulation, while descriptions of extended spatial configurations require the selection of coordinate systems or spacelike hypersurfaces.

A foliation divides a spacetime model into hypersurfaces assigned successive values of a time parameter. Such a construction is mathematically useful. It is not, merely by being constructed, the uniquely privileged physical division of the universe into successive universal moments.

Certain cosmological solutions possess a natural cosmic-time parameter associated with their symmetry and a specified family of observers. This supplies a useful coordinate structure for cosmology. It does not establish an absolute universal present as a general requirement of relativistic physics.

The distinction between mathematical representation and physical ontology is essential.

\begin{proposition}[No privileged universal present]

Special relativity does not supply an observer-independent simultaneity relation for all spatially separated events. General relativity does not, through its field equations alone, restore a universally preferred foliation of spacetime into physical presents.

\end{proposition}

\begin{proof}

The Lorentz transformation demonstrates that distant simultaneity depends upon inertial reference frame. Therefore no frame-independent assignment of distant simultaneous events is supplied by special relativity. General relativity describes physical relationships through a spacetime metric and permits coordinate and foliation choices without generally designating one such choice as a privileged universal present. Consequently, the relativistic laws do not require a single universe-wide physical now.

\end{proof}

This proposition establishes the foundation of the earlier McKinley paper \emph{The Present Was Never Universal} \cite{PresentNeverUniversal}.

Its significance for the present investigation is substantial.

If there is no privileged universal present, there is no privileged universal succession of presents obtained simply by advancing that present along a global clock.

An observer may construct a mathematical timeline for the universe. Another observer may construct a different foliation. The resulting descriptions remain related by the governing relativistic structure.

Neither construction becomes an actual collection of universal physical moments merely because the mathematics represents it.

The familiar filmstrip picture has lost its physical foundation.

\section{Physical Actuality Is Local}

The statement that physical actuality is local requires explanation. Without such explanation, the phrase risks becoming an abstract slogan rather than a physical claim.

Begin with an ordinary clock.

The clock occupies a physical trajectory through spacetime. Its internal condition changes according to the interactions and physical laws governing its operation. Its displayed reading reflects accumulated proper time along that trajectory.

Another clock occupies another trajectory. Its accumulated proper time reflects its own path and gravitational circumstances. The clocks can be compared when their readings become available within a common causal interaction, but no comparison makes their preceding proper-time accumulations identical.

The same principle applies more broadly to physical systems.

At any particular physical location, the locally available configuration reflects the lawful physical interactions associated with that location. Material conditions, electromagnetic interactions, gravitational circumstances, incoming information, and retained physical records contribute to what is locally registered.

Another location possesses different causal circumstances. It is reached by different signals, occupies a different geometrical position, and participates in a different collection of local interactions.

The two locations do not require identical physical information or identical temporal circumstances.

Strictly speaking, an abstract spatial point does not itself carry a clock or accumulate proper time. Proper time belongs to timelike worldlines followed by physical systems. The local physical configuration is therefore understood through the systems and interactions present in the relevant spacetime neighborhood.

This is the physically precise meaning of local temporal accumulation.

The current condition of a material system reflects its lawful evolution along its worldline and the causal interactions in which it has participated. Its physical history is not an amount of universal time deposited into the system. It is the set of lawful relationships through which its present condition has arisen.

The term \emph{micro-local} is useful for emphasizing that locality begins at the smallest physical scale relevant to the description, rather than at the scale of a human observer or laboratory instrument. The argument does not require a particular microscopic spacing of spacetime or a new quantum-gravity hypothesis. The essential principle is that physical conditions and interactions are locally instantiated and causally organized.

An extended detector consists of physically coupled constituents. Its output integrates many local interactions into one coordinated record. A human observer performs the same kind of integration on a vastly more complicated scale.

The resulting experience of a continuous world is real as a lawful organization of physical information. It does not require every distant location in the universe to share the observer's moment.

Consider an observer looking through a telescope.

The observer registers a present physical signal at the detector. The registered information is lawfully related to conditions associated with a distant astronomical source. The observer interprets that information using the causal structure of spacetime and the relevant physical laws.

The distant source is not thereby brought into the observer's local present as an entire simultaneous physical configuration. What is locally registered is the information available through the lawful relation.

This point was developed in the earlier McKinley registration papers, where the invariant speed \(c\) was interpreted as the causal bound governing observer-side information availability \cite{RegistrationBound}.

The present paper extends the same principle to physical history.

\begin{definition}[Local physical actuality]

Local physical actuality is the physically instantiated condition of a system or interaction at a spacetime location, including its locally available records and causal relationships. Under the present-only interpretation, such actual conditions are not treated as portions of one privileged universe-wide instant.

\end{definition}

This definition does not make physical actuality dependent upon human consciousness. A detector, atom, material interaction, or other physical system participates in registration regardless of whether a person observes it.

Nor does it mean that every location operates according to different physical laws.

The laws remain coherent. What differs is the locally instantiated physical condition, the available information, and the temporal and causal relationships associated with each physical system.

A useful distinction therefore emerges.

The lawful structure governing the universe is common. The local registrations governed by that structure are distinct.

The present is not one enormous physical sheet extending across the cosmos. In the proposed ontology, physical actuality is locally instantiated, and the coherent relationships among those local conditions constitute the observable physical world.

\section{No-Go on a Global Sequence of Physical Presents}

The absence of universal proper-time accumulation and privileged universal simultaneity has an immediate consequence for the familiar concept of a global physical history.

A conventional dynamical model represents a system through a state function,

\begin{equation}
S(t).
\end{equation}

The function associates a configuration with each value of the time parameter. By evaluating the function at different times, the physicist obtains a sequence,

\begin{equation}
S(t_0),\quad S(t_1),\quad S(t_2),\quad \ldots
\end{equation}

For an ordinary laboratory system, such a description is convenient and highly effective.

The difficulty arises when the entire universe is imagined as possessing one privileged sequence of configurations of this kind.

A universal configuration at one instant requires a criterion assigning spatially separated events to the same temporal slice. Relativity does not supply one uniquely privileged criterion.

A relativistic spacetime model may be foliated into hypersurfaces \(\Sigma_t\). The physical fields and other relevant quantities may then be represented on those hypersurfaces. Different admissible foliations describe the same underlying relativistic relationships through different assignments of global coordinate time.

The existence of these descriptions does not identify one foliation as a sequence of uniquely actual universal moments.

\begin{proposition}[No privileged global succession]

Relativistic physics does not require one privileged sequence of simultaneous configurations representing the successive universal presents of the cosmos.

\end{proposition}

\begin{proof}

A privileged sequence of universal presents requires a privileged assignment of distant simultaneity at each step of the sequence. By the preceding relativistic result, no such assignment is supplied generally by special or general relativity. A sequence of spacelike hypersurfaces is therefore a mathematical organization of the spacetime description, not an invariantly privileged succession of universe-wide presents.

\end{proof}

This is a no-go on a particular physical ontology, not on the existence of causal history.

Relativity still describes timelike and null relationships among events. It distinguishes causal precedence from spacelike separation. It provides invariant intervals, proper-time durations where applicable, and lawful constraints on information propagation.

What it does not provide is a universal movie advancing one complete cosmic frame at a time.

The distinction matters because the ordinary picture of a global past begins by assuming that the universe once possessed one complete universal present and then moved away from it.

If the universal present was never established, the supposed succession of former universal presents has no independent foundation.

The past cannot be identified with a completed stack of universal moments when the underlying universal moments were never physically privileged in the first place.

\section{The Timeline Is a Calculation}

The absence of a privileged global present does not prevent physics from constructing timelines.

A physicist describes a satellite trajectory by assigning positions and velocities to successive values of a coordinate time. An astronomer reconstructs an earlier planetary arrangement by applying the relevant dynamical equations. A cosmologist employs a model of large-scale expansion to calculate earlier conditions of the observable universe.

These calculations are not defective because the universe lacks a universal clock. They are descriptions formulated under specified coordinates, assumptions, and physical laws.

The timeline is a mathematical instrument.

Its existence as a representation does not establish that each represented configuration is physically present somewhere.

A useful analogy is a graph of temperature.

Suppose a laboratory instrument records the temperature of a vessel throughout an experiment. A graph shows the temperature at noon, one o'clock, two o'clock, and three o'clock.

The graph remains available at four o'clock. The vessel's earlier temperature conditions do not thereby remain physically present inside the graph.

The curve identifies the relationship between temperature and assigned time. It does not preserve the earlier physical state of the vessel.

A physical timeline performs the same descriptive function on a broader scale.

A complete mathematical model may assign quantities to every event in a spacetime manifold. This permits the calculation of intervals, causal relationships, trajectories, and field conditions throughout the modeled domain.

But a mathematical assignment is not itself an additional physical registration.

The model describes what the theory assigns to the events. Whether all represented events possess equal and continuing physical existence is a separate ontological question.

The present paper rejects the automatic promotion of a mathematical timeline into a physical archive.

A reconstruction of yesterday is not yesterday.

A prediction of tomorrow is not tomorrow.

The ability to calculate a physical configuration does not cause that configuration to exist.

This distinction is already familiar in other areas of physics. A mathematical formalism describes permitted or inferred physical relations without every mathematical object in the formalism becoming an independently existing physical object.

The same discipline should be applied to time.

\section{What Physics Actually Possesses of the Past}

If the past is not an existing location, what does physics possess when it investigates earlier events?

It possesses physical evidence.

Consider a vase that stood intact on a table yesterday and lies broken on the floor today.

An investigator enters the room and examines the fragments. Their shapes, locations, and material properties supply evidence concerning the earlier configuration. A photograph might show the vase before it broke. A camera recording might show its fall. A witness might remember what happened.

All of these are physical conditions available to the investigator.

The fragments are present physical arrangements. The photograph is a present arrangement of recorded information. The camera file is a present information-bearing state of a storage system. The witness's memory is represented through presently existing physical processes and structures.

None of these is the intact vase as it existed yesterday.

The investigator does not recover yesterday's physical configuration as an object. The investigator reconstructs it from surviving evidence.

The same distinction applies to a fossil.

The fossil is a present physical structure whose characteristics support conclusions about an earlier organism. It does not contain the earlier organism's complete living physical configuration.

The same distinction applies to astronomical observation.

A detector registers physical information lawfully associated with conditions at a distant source. The corresponding emission conditions are inferred through the governing physical relationships. The observation does not place the observer in an earlier physical universe.

A clock likewise provides a current physical indication. Its reading expresses accumulated proper time according to its trajectory and the laws governing its operation. The clock does not contain the hours as physical objects.

These examples reveal a general principle.

\begin{proposition}[Historical evidence is presently registered information]

The physical evidence available to an observer concerning an earlier event consists of records, signals, material consequences, and other information-bearing physical conditions registered at the observer's location or otherwise made causally accessible to that observer.

\end{proposition}

\begin{proof}

Physical inspection requires interaction with physically available conditions. A photograph is inspected through present optical or electronic registration. A fossil is inspected through present material interactions. A memory is accessed through present physical processes. An astronomical observation occurs through a present detector registration. In each case, the earlier event is inferred from the available information rather than directly inspected as an earlier event newly occupying the observer's location.

\end{proof}

The consequence is not that physical history is fictitious.

It is that historical knowledge and the continued physical existence of historical configurations are different claims.

Physics possesses information from which it reconstructs what occurred. That information does not constitute the earlier occurrence itself.

\section{Historical Reconstruction and Its Information Limits}

Historical reconstruction applies physical law to available evidence in order to infer earlier conditions.

Let \(H\) denote a candidate physical history, and let \(R(H)\) denote the information from that history preserved in the records available to an observer.

The reconstruction problem is to infer \(H\) from \(R(H)\) under the applicable physical laws.

An important limitation follows when different candidate histories are compatible with the same surviving information.

Suppose

\begin{equation}
H_1\neq H_2
\end{equation}

while

\begin{equation}
R(H_1)=R(H_2).
\end{equation}

The available records do not distinguish the histories.

This is not a claim that the governing physics must be fundamentally irreversible, nor that complete microscopic information must always be lost. The point concerns the information actually retained and available for reconstruction.

A simple example is a computer storage device that has been repeatedly overwritten.

Suppose one set of files occupied the device, was deleted, and was replaced by another. This happened repeatedly. Eventually, imagine that no accessible physical distinction remains that identifies the earlier contents.

An investigator who receives the final device examines its current physical configuration.

The earlier sequences of files are not recovered merely because the investigator knows that overwriting occurred.

Different sequences could be compatible with the accessible final record. Without distinguishing evidence, the investigator cannot identify the unique earlier sequence.

A similar example involves a ball moving through a pachinko machine.

The ball encounters pins, changes direction repeatedly, and eventually arrives at a particular location. An observer sees only its final position.

The final position does not necessarily determine the exact earlier route. Different sequences of collisions can lead to the same observed destination.

The investigator constructs possible histories. The final position does not itself contain the complete preceding trajectory.

The physical universe presents a more complicated reconstruction problem, but the logical relationship is the same.

A present record supplies information. A model supplies lawful relationships. The resulting historical reconstruction depends on both.

\begin{proposition}[No unique reconstruction from insufficient records]

If two distinct physically admissible histories produce identical information in the surviving records available to an observer, those records do not uniquely determine which history occurred.

\end{proposition}

\begin{proof}

Let \(H_1\) and \(H_2\) be distinct histories compatible with the governing laws, and suppose \(R(H_1)=R(H_2)\). An inference based solely upon \(R(H)\) and the laws common to both histories receives identical inputs in both cases. It therefore cannot select one history uniquely without additional distinguishing information or assumptions. The available record is insufficient for unique reconstruction.

\end{proof}

This establishes an informational no-go.

It does not establish that earlier events did not occur. It establishes that the information available now need not identify every earlier physical configuration.

The distinction is especially important when discussing the earliest universe.

The existence of the current universe does not, by itself, supply a complete record of every earlier condition. Nor does the ability to construct a successful cosmological model automatically establish that the model identifies an absolutely first physical event.

The amount of history that physics reconstructs depends upon the evidence and upon the lawful relationships governing its interpretation.

But there is a still more fundamental question.

Even if the earlier history were reconstructed perfectly, would that establish that the earlier physical configurations continue to exist?

The answer is no.

\section{The Central No-Go: Reconstruction Does Not Establish an Existing Past}

The central result of this paper concerns the difference between reconstructing a physical configuration and establishing its present physical existence.

Suppose a physicist possesses sufficient information to calculate the exact physical configuration of a laboratory system at an earlier time.

Let the reconstructed configuration be denoted by

\begin{equation}
S(t_0).
\end{equation}

Let the physically available evidence be \(R\), and let the governing lawful structure be \(L\).

The reconstruction may be represented schematically as

\begin{equation}
(R,L)
\longrightarrow
S(t_0).
\end{equation}

The expression means that the evidence and physical laws support the reconstruction of a configuration assigned to an earlier time.

It does not mean that the reconstruction has physically created or recovered that earlier configuration.

Nor does it establish that the earlier configuration was retained somewhere throughout the intervening duration.

The distinction becomes obvious when applied to ordinary objects.

A physicist can calculate where a projectile was located ten seconds ago. The calculation does not establish that the projectile still occupies that earlier location as a second physical object.

An engineer can reconstruct the configuration of a machine before it failed. The reconstruction does not restore the earlier functioning machine to physical existence.

An astronomer can calculate the position of a planet at an earlier date. The calculation does not establish that the earlier planet and its surroundings remain somewhere as a presently existing configuration.

In each case, the calculation produces a description.

The earlier condition is the subject of the description.

It is not a second physical thing that appears because the description is accurate.

This can be stated formally.

\begin{proposition}[No-go on persistent-past inference]

The successful reconstruction of an earlier physical configuration from currently available evidence and governing physical laws does not entail that the earlier configuration continues to exist anywhere in physical reality.

\end{proposition}

\begin{proof}

Historical reconstruction establishes a relationship between available physical evidence, governing laws, and an inferred earlier configuration. The evidence consists of presently accessible physical information. The reconstruction is a mathematical or descriptive result concerning conditions assigned to an earlier event or interval.

The proposition that the inferred configuration continues to exist is a further ontological proposition. It does not follow from the accuracy of the mathematical description or from the existence of the current record.

A present-only interpretation and an interpretation assigning continuing existence to all events of a spacetime description can employ the same relativistic equations and reconstruct the same earlier configurations. The reconstruction operation alone therefore does not distinguish those ontologies or entail continuing existence.

Consequently, successful reconstruction does not establish the physical persistence of the reconstructed earlier configuration.

\end{proof}

The no-go is directed against a specific inference.

It does not derive the nonexistence of earlier configurations solely from the mathematics of relativity. It establishes that the mathematics and the surviving evidence do not require their continued existence.

The proposed present-only ontology takes the further position that earlier configurations do not remain physically actual.

This is a deliberate distinction between an established inferential limit and the physical interpretation advanced here.

The resulting picture is economical.

Physics requires lawful relationships among physical conditions. It requires causal structure and proper-time accumulation where applicable. It requires records from which earlier events are inferred and laws under which future possibilities are calculated.

It does not require a stored collection of previous physical configurations.

The previous configuration is not hidden inside the current configuration. The current configuration contains whatever physical consequences and records have survived.

The past is not waiting somewhere behind the present.

There is no physical archive merely because a mathematical timeline can be drawn.

\section{No Global Past}

A further conclusion follows from the combined relativistic and historical arguments.

The ordinary phrase \emph{the history of the universe} encourages the idea of one universal sequence through which the cosmos has progressed.

Yet there is no universally accumulated proper time, and there is no privileged relativistic simultaneity assigning one physical moment to every location.

The physical universe therefore does not possess an invariantly privileged succession of complete simultaneous configurations.

There are local physical relationships, timelike worldlines, null relationships, and causal orderings governed by relativistic structure.

The term \emph{past} remains meaningful within those relationships. An event can lie in the causal past of another event. A clock can record a duration since an earlier interaction. A fossil can preserve information concerning an earlier organism. None of these uses requires a universally existing physical past.

The causal past of a detector event is defined by the spacetime events capable of influencing that detector event through timelike or null relations. It is a geometrical and causal relation within the model.

It is not a warehouse containing earlier versions of physical reality.

Similarly, a full spacetime solution provides a mathematical description of event relationships across a modeled domain. The solution need not be interpreted as a literal collection of coexisting physical moments.

The local-present model takes this separation seriously.

At each locally actual registration, physical conditions possess definite lawful relationships. Some of those conditions are records whose contents refer to preceding events. The content of the record does not make the preceding event current.

The phrase \emph{no global past} therefore has two closely related meanings in this paper.

First, relativity does not supply a privileged universal sequence of earlier simultaneous configurations.

Second, the present-only ontology rejects the additional picture of the past as a physically existing collection of configurations located somewhere along that sequence.

The second claim is not the definition of the first. It is the physical interpretation developed from the first together with the record-reconstruction distinction.

This separation preserves the strength of the argument.

Relativity establishes the absence of a privileged universal temporal slicing. Historical reconstruction establishes lawful relationships between surviving evidence and inferred earlier events. Neither result requires the persistence of the past.

The proposed physical model therefore assigns no location to a physically existing past.

The past is not somewhere else.

It is not somewhere at all.

\section{There Is No Universal Clock from Which the Photon Needs an Exemption}

The foregoing argument has an immediate consequence for the interpretation of light.

The familiar picture of a photon begins with a small physical object traveling through space. The object departs an emitter, moves across an intervening distance, and eventually arrives at an absorber.

Observers measure a temporal separation between emission and absorption. The separation may be extremely short or extraordinarily long.

Under the ordinary universal-clock picture, the photon seems to travel through the universe's elapsed time while somehow experiencing none of it.

This produces a familiar conceptual puzzle.

How does something cross an enormous distance without accumulating time?

The puzzle arises because the question begins by assigning universal temporal progression to everything in physical reality.

Relativity has already removed that assignment.

Proper time is accumulated along timelike worldlines. Different time-bearing systems accumulate different proper times. Null worldlines possess zero proper-time interval.

For a null relation in flat spacetime,

\begin{equation}
ds^2
=
-c^2dt^2+dx^2+dy^2+dz^2
=
0.
\end{equation}

The corresponding proper-time increment is

\begin{equation}
d\tau=0.
\end{equation}

Thus, along a null curve,

\begin{equation}
\Delta\tau_{\mathrm{null}}=0.
\end{equation}

This is the standard relativistic result.

A photon possesses no rest frame. It accumulates no proper time along a null trajectory. The ordinary notion of a photon experiencing a sequence of internal moments is therefore not supplied by the relativistic description.

The Timeless Light Model advances the further interpretation that the photon is not a persisting traveler undergoing successive internal states between emission and absorption. It is a lawful relation whose emitter-side and absorber-side state changes are registered within spacetime \cite{NullProperTime,UnifiedPhysics,EndpointsTimed}.

The distinction between the registered events and the relation connecting them is central.

An emitter undergoes a state change associated with emission. An absorber undergoes a state change associated with absorption. These events possess spacetime coordinates and lawful causal relationships.

An observer can assign an elapsed coordinate time between them.

That elapsed time is a feature of the spacetime description. It is not proper time accumulated by the photon.

The invariant limit \(c\) governs the null structure and the causal availability of information at registered locations.

It does not supply a hidden universal clock through which the photon must progress.

The new paper's contribution to this interpretation is to remove the mental picture that made the null result appear exceptional.

If the entire universe were progressing through one universal sequence of moments, the photon would seem to require an exemption from that progression.

But there is no such universal progression established by relativity.

Time-bearing systems possess their own accumulated proper times. Null relations do not.

The photon is not an exception to a universal clock.

There is no universal clock from which the photon needs an exemption.

\begin{proposition}[Null proper time requires no universal-clock exception]

The zero proper-time interval associated with a null relation does not conflict with the existence of accumulated proper time along timelike worldlines, because relativity assigns proper time according to worldline structure rather than according to a universal elapsed duration.

\end{proposition}

\begin{proof}

Proper time is defined by the metric interval along a worldline. For timelike worldlines, its integral is positive over a nonzero timelike segment. For null curves, the metric interval vanishes, so the accumulated proper time is zero. Since no universal proper-time accumulation is required for all worldlines, the null result does not violate such a requirement. The apparent exception arises only when universal temporal accumulation is imposed as an additional assumption.

\end{proof}

This proposition does not, by itself, prove the full Timeless Light Model interpretation of photon ontology. The reclassification of the photon as a lawful relation rather than a persistent traveler remains the interpretive contribution of the earlier corpus.

What the proposition establishes is that the absence of photon-side proper time is not anomalous within the relativistic architecture.

The wider present-only model reinforces that result by removing the image of a universally advancing physical filmstrip.

The emitter and absorber possess registered spacetime relationships. The photon relation has no accumulated internal time. Nothing in the physical description requires an additional universal clock against which that absence must be explained.

\section{Lawful Continuity Does Not Require a Stored Past}

A natural question arises once the physical past is removed from the proposed ontology.

How does the universe preserve continuity?

A stone remains a stone. A planet continues in its orbit. A radioactive atom has a lawful decay structure. A biological organism retains organized physical relationships. A computer preserves information through physical storage states.

These familiar continuities might appear to require earlier physical configurations to remain somewhere in existence.

They do not.

A material system possesses a presently actual physical condition governed by physical law. Its lawful description includes relationships among states assigned to different events. Physical records and stable material structures preserve information relevant to those relationships.

The governing law does not need to retrieve an earlier physical object from a stored past in order to describe the subsequent state.

Consider a conventional evolution equation,

\begin{equation}
\frac{dS}{dt}=F(S,t).
\end{equation}

The equation specifies a lawful relationship among states described at different parameter values. Given suitable information and conditions, the physicist uses it to reconstruct earlier values or calculate subsequent values.

The equation does not state that all those values are independently existing physical configurations.

Its mathematical function is to represent the lawful relation.

The same distinction applies to relativistic field equations and quantum dynamics.

A spacetime description identifies lawful relationships among events and fields. A quantum description assigns amplitudes and probabilities to lawfully admissible outcomes. A local detector registers a particular physical result.

The mathematical structures are not themselves proof that every represented event remains physically actual.

This is consistent with the clarified architecture developed in \emph{Physics Has Always Been Unified} \cite{UnifiedPhysics}.

In that architecture, lawful structure governs physical admissibility. Quantum mechanics describes lawful eligibility and the amplitude and probability relationships of admissible outcomes. Relativity governs the spacetime organization of registration.

The governing laws are not physical objects aging through a universal time.

Nor does a law need to preserve an earlier configuration as an existing object in order to govern lawful physical continuity.

A stable atom retains its lawful structure through the physical relationships constituting its condition. A clock accumulates its own proper time. A computer stores information in its current physical configuration.

None requires yesterday to remain physically present somewhere.

The distinction is especially important for memory.

Memory is a present physical information-bearing condition related to prior events. Remembering is a present physical activity. The remembered event does not return merely because the memory becomes active.

The same applies to all historical information.

The universe does not need to retain its earlier configurations as physical things. It needs only the physical conditions and lawful relationships sufficient for the processes and observations under examination.

The present-only interpretation thus removes a presumed additional layer of reality without changing the governing mathematics.

\section{Cosmology and the Question of a Beginning}

The distinction between physical history and surviving evidence becomes particularly important in cosmology.

Modern cosmology reconstructs an earlier hot, dense phase of the observable universe. The evidence includes large-scale cosmic expansion, the cosmic microwave background, and primordial elemental abundances.

These observations support detailed models of early physical conditions.

The success of those reconstructions is not challenged here.

The question is what the observations establish concerning physical existence.

A cosmological model assigns physical quantities to earlier epochs. It may identify conditions of extreme density and temperature. It may describe the evolution of the scale factor and the relationships among matter, radiation, and spacetime geometry.

The model is a mathematical and evidential reconstruction.

It does not thereby demonstrate that the earlier cosmic configuration remains physically existent somewhere.

Nor does reconstruction of an early hot, dense phase, by itself, settle whether physical existence possessed an absolute beginning.

The question of the first configuration and the question of the earliest reconstructible configuration are not identical.

Suppose the available evidence supports reconstruction back to some modeled condition \(S_0\).

That condition is identified through the model and the information available for inference.

Whether \(S_0\) was absolutely the first physical condition is a further question. Whether \(S_0\) remains physically existent is another.

The present paper addresses the latter inference directly and separates it from the former.

One can reconstruct an event without physically revisiting it. One can identify an earlier condition without establishing that it survives. One can model a beginning without turning the beginning into an existing location.

This distinction applies equally to models possessing an initial temporal boundary and models extending indefinitely into earlier parameter values.

The present-only interpretation does not require the universe to possess a complete recoverable physical history. It does not require a decision concerning an absolute beginning in order to identify the distinction between present evidence and historical reconstruction.

The central result remains unchanged.

Cosmological observations are physically registered evidence. Cosmological models reconstruct lawful earlier conditions. Those reconstructions do not establish that the earlier universe still exists somewhere.

\section{What the No-Go Results Establish}

The no-go arguments developed here address distinct assumptions, and their logical relationships should be made explicit.

The first result concerns proper time. Relativity positively establishes that time-bearing systems following different timelike worldlines can accumulate different durations. It therefore supplies no universal accumulation of proper time shared by every physical clock.

The second result concerns simultaneity. Relativity supplies no invariantly privileged distant simultaneity and no generally preferred global foliation. Consequently, it does not require one universal sequence of simultaneous physical configurations.

The third result concerns information. Where different candidate histories produce the same accessible records, those records do not uniquely identify which history occurred. The result establishes a conditional limit on reconstruction from insufficient information.

The fourth result concerns ontology. Even successful reconstruction does not establish that the reconstructed earlier configuration continues to exist anywhere in physical reality.

The logical strength of the fourth no-go lies in the distinction between description and existence.

A mathematical model can represent an earlier event. A physical record can support the conclusion that the event occurred. Neither fact entails that the event remains physically actual.

The full present-only ontology goes beyond this inferential result by affirmatively rejecting a still-existing physical past.

That interpretive commitment is not presented as a direct mathematical theorem of special or general relativity. Relativistic equations admit four-dimensional spacetime descriptions, and those mathematical descriptions do not, on their own, settle the ontological status of every represented event.

The present paper therefore distinguishes the physical results from the interpretation selected.

The physical results remove the necessity of a universal clock, a privileged global present, and a physically preserved past.

The proposed interpretation then adopts the economical reading that physical actuality consists of locally instantiated conditions governed by timeless lawful structure.

This reading does not require a second universe, a hidden temporal storage region, or an additional mechanism for preserving vanished physical configurations.

There is no physical past to store.

There is physical evidence from which earlier configurations are reconstructed.

The distinction is simple, but its consequences are substantial.

\section{Discussion: Removing the Filmstrip from Physics}

The argument developed here is principally an exercise in removing an unnecessary physical picture.

The mathematical description of the universe is not being discarded. Relativity continues to govern spacetime geometry, proper time, and causal order. Quantum mechanics continues to describe the lawful eligibility of physical outcomes and their mathematical probability relationships. Historical sciences continue to reconstruct earlier events from surviving records.

What is removed is the assumption that these descriptions require the entire represented timeline to exist physically.

The distinction resembles other cases in which a successful mathematical representation acquires an unwarranted physical interpretation.

A graph is not the process it depicts. A map is not the territory it represents. A reconstruction is not the configuration reconstructed.

A timeline is not an archive of existing moments.

The ordinary universal-clock picture makes the universe appear to be a film advancing through successive frames. The observer imagines the current frame as the present, the completed frames as the past, and the unshown frames as the future.

Relativity provides no preferred universal frame of this kind.

Instead, proper times belong to timelike worldlines. Causal relationships govern which information becomes available at different locations. Physical registration occurs locally, and its results are coordinated by common lawful structure.

The temporal description of one location is therefore not simply a small portion of a universal clock shared by every other location.

The presently registered conditions at different physical systems reflect their own worldlines, interactions, and available information.

The universal filmstrip is unnecessary.

The resulting model is not a denial of physical reality outside human observation. Registration is understood physically, not psychologically. An interacting atom, a material system, or a detector participates in local physical relationships whether or not a person witnesses the result.

Nor is the model a rejection of lawful historical relationships.

An earlier event remains earlier in the relevant causal or proper-time ordering. A clock's recorded duration is meaningful. The reconstruction of an astronomical event is meaningful. The causal ancestry of an observed structure is meaningful.

What is denied is that an earlier event must remain physically existent in order for its causal or mathematical relationships to be meaningful.

The distinction also helps clarify the proper-time foundation of the Timeless Light Model.

The photon appears anomalous only when every physical relation is assumed to move through one universal clock. Once that clock is removed, null proper time is simply the appropriate relativistic temporal structure for a null relation.

The remaining interpretive question concerns what kind of physical relation the photon represents.

The Timeless Light Model answers that question through lawful structure and spacetime registration. The present paper supplies the broader temporal picture in which that answer becomes easier to understand.

There is no universal accumulation of time to which the photon must conform.

There is no physically privileged universal present through which it must pass.

There is no necessary collection of earlier universal moments remaining in existence behind its registered absorption event.

There is lawful structure, relativistic causal order, and the registered physical conditions governed by them.

The present-only model therefore replaces the image of the universe progressing through an existing timeline with a more economical description: local physical actuality organized by timeless law.

\section{Conclusion}

The universe is commonly imagined as possessing a past, a present, and a future arranged along a timeline. The picture is useful for calculation and historical description, but its apparent simplicity conceals several physical assumptions.

Relativity does not supply one accumulated proper time shared by all physical systems. Proper time belongs to timelike worldlines, and different worldlines can accumulate different durations.

Relativity also does not supply a privileged universal present. Spatially separated simultaneity depends upon reference frame, and the mathematical assignment of spacelike hypersurfaces does not establish one universally privileged succession of cosmic configurations.

Physical registration is local. The condition of a time-bearing system reflects its lawful spacetime history and causal interactions. Information becomes available to physical systems according to the applicable causal structure, including the invariant limit \(c\).

Physics uses presently available records and governing laws to reconstruct earlier physical configurations. It uses the same lawful structure to calculate subsequent possibilities.

These operations do not establish that the configurations being reconstructed or predicted exist as additional physical realities.

The paper has therefore established a no-go on the inference from successful historical reconstruction to a continuing physical existence of the past.

The proposed present-only ontology takes the further step of removing the past from the inventory of physically actual things.

Earlier events occurred in their proper causal and temporal relationships. Their physical consequences and surviving records contribute to conditions available now. But the earlier configurations are not preserved as existing objects occupying another location along a timeline.

The same interpretation illuminates the photon.

A photon relation accumulates no proper time. This is not a mysterious exemption from a universal temporal progression. The supposed universal progression was never supplied by relativity.

There is no universal clock from which the photon needs an exemption.

The universe does not require its past to remain physically existent in order to possess continuity, lawful structure, causal relationships, historical evidence, or reconstructible history.

The universe is not a filmstrip with every frame still available.

It is a coherent physical order in which actuality is locally instantiated, temporal accumulation belongs to time-bearing systems, and causal information is registered according to relativistic law.

The past is reconstructed from physical evidence.

The future is calculated from physical law.

Neither is a second layer of present physical actuality.

The past is not somewhere else.

\emph{It is not somewhere at all.}

\section*{Relationship to Prior McKinley Results}

The present paper extends, rather than replaces, the earlier Timeless Light Model results.

\emph{The Present Was Never Universal} established the absence of a privileged universe-wide present and developed the interpretation of local physical registration. The present paper uses that relativistic foundation to examine the status of global historical sequences and the physical meaning of the past.

\emph{Physics Has Always Been Unified} established the clarified architecture in which timeless lawful structure governs physical admissibility, quantum mechanics describes lawful eligibility, and relativity governs the spacetime organization of registered outcomes. The present paper applies that distinction to historical reconstruction and the ontology assigned to earlier physical configurations.

The previously developed null-proper-time and photon-ontology papers provide the related interpretation of light as a lawful relation rather than an internally time-bearing traveler. The present paper does not rederive that full interpretation. It shows why abandoning a universal-clock picture makes the null proper-time result conceptually natural.

No new physical constant, interaction, dynamical equation, or experimental prediction is proposed. The contribution is an interpretive clarification of time, physical actuality, and historical reconstruction under the existing laws.

\begin{thebibliography}{99}

\bibitem{Einstein1905}
A.~Einstein,
``Zur Elektrodynamik bewegter K\"orper,''
\emph{Annalen der Physik} \textbf{17},
891--921 (1905).

\bibitem{Minkowski1909}
H.~Minkowski,
``Raum und Zeit,''
\emph{Physikalische Zeitschrift} \textbf{10},
75--88 (1909).

\bibitem{Rindler2006}
W.~Rindler,
\emph{Relativity: Special, General, and Cosmological},
2nd ed.,
Oxford University Press (2006).

\bibitem{Wald1984}
R.~M.~Wald,
\emph{General Relativity},
University of Chicago Press (1984).

\bibitem{MTW1973}
C.~W.~Misner, K.~S.~Thorne, and J.~A.~Wheeler,
\emph{Gravitation},
W.~H.~Freeman (1973).

\bibitem{Malament2012}
D.~B.~Malament,
\emph{Topics in the Foundations of General Relativity and Newtonian Gravitation Theory},
University of Chicago Press (2012).

\bibitem{PresentNeverUniversal}
J.~C.~W.~McKinley,
\emph{The Present Was Never Universal:
Local Registration and the Absence of a Universe-Wide Now},
manuscript, July 2026.

\bibitem{UnifiedPhysics}
J.~C.~W.~McKinley,
\emph{Physics Has Always Been Unified:
General Timelessness, Null Proper Time,
and the Clarified Architecture---Why Quantum Mechanics
and Relativity Were Never in Competition},
Zenodo (2026).
\doi{10.5281/zenodo.20954931}.

\bibitem{GeneralTimelessness}
J.~C.~W.~McKinley,
\emph{Timelessness Is the General Condition:
Spacetime Is the Special Case},
Zenodo (2026).
\doi{10.5281/zenodo.19771924}.

\bibitem{NullProperTime}
J.~C.~W.~McKinley,
\emph{Taking Null Proper Time Seriously:
An Interpretive Clarification of Null Proper Time},
Zenodo (2025).
\doi{10.5281/zenodo.18004632}.

\bibitem{RegistrationBound}
J.~C.~W.~McKinley,
\emph{c Is a Registration Bound, Not a Traveler's Speed:
A Registration-Based Interpretation of the Information Speed Limit},
Zenodo (2026).
\doi{10.5281/zenodo.20114176}.

\bibitem{EndpointsTimed}
J.~C.~W.~McKinley,
\emph{The Endpoints Are Timed; The Relation Is Not},
working manuscript, September 2026.

\end{thebibliography}

\end{document}
